En los últimos años, la impresión 3D ha pasado de ser una tecnología exclusiva de laboratorios especializados a convertirse en una herramienta presente en entornos industriales, educativos y domésticos. Esta popularidad ha generado una amplia oferta de arquitecturas comerciales orientadas al consumo general.

Sin embargo, gran parte de las impresoras 3D comerciales disponibles actualmente están orientadas a usuarios no especializados y priorizan una operación inmediata con mínima intervención técnica. Como consecuencia, muchos de estos equipos utilizan arquitecturas cerradas, firmware restringido y componentes propietarios que limitan las posibilidades de modificación y reparación del sistema por parte del usuario.

En entornos académicos y de investigación, esta dependencia tecnológica genera diversas problemáticas. La necesidad de utilizar refacciones específicas del fabricante incrementa los costos de mantenimiento y dificulta la reparación del equipo ante fallas menores. De igual manera, los tiempos asociados a procesos administrativos de adquisición pueden provocar largos periodos de inactividad en laboratorios universitarios. Desde el punto de vista formativo, el uso de sistemas cerrados limita la capacidad del estudiante para comprender e intervenir directamente en aspectos mecánicos, electrónicos y de control relacionados con la manufactura aditiva.

Frente a esta problemática, las plataformas basadas en hardware y software de código abierto representan una alternativa viable para el desarrollo de sistemas modificables, reparables y de bajo costo. La posibilidad de fabricar componentes propios y adaptar tanto el firmware como la arquitectura mecánica permite extender la vida útil del equipo y fortalecer el aprendizaje práctico en áreas relacionadas con control numérico computarizado (CNC), electrónica y diseño mecánico. Esta alternativa es especialmente pertinente en instituciones con recursos limitados, pues permite al estudiante interactuar con cada subsistema del equipo, diagnosticar fallas y proponer mejoras documentadas.

Por lo tanto, surge la necesidad de desarrollar una impresora 3D tipo FDM con arquitectura abierta y fácilmente mantenible, orientada a aplicaciones académicas, capaz de servir como plataforma didáctica para la experimentación y el aprendizaje práctico. Abordar esta necesidad es la parte central del presente proyecto, que busca no solo la construcción de un equipo funcional, sino también la creación de un entorno didáctico que permita a los estudiantes adquirir competencias en diseño, manufactura y control de sistemas de manufactura aditiva.


\begin{comment}
La impresión 3D, en todas sus presentaciones, emergió como una de las 
tecnologías mas emocionantes en la actualidad, mejorando el alcance de la 
fabricación y creación de objetos de manera automatizada. Con origen en 
laboratorios de investigación, llegó a convertirse en una 
herramienta omnipresente en industria e incluso en talleres y hogares de 
aficionados a la fabricación de todo el mundo. Debido a su aumento de 
popularidad, su oferta en el mercado aumentó, por lo cual vemos 
arquitecturas comerciales al alcance de cualquiera, conformándose como un
producto no solo para un sector sino que se proyecta como una maquina que
puede llegar a ser común en los hogares.

¿Cual es el problema?
El hecho de que estas maquinas están destinadas a usuarios un tanto mas 
casuales, como lo seria un investigador en un laboratorio o a personal 
de una fabrica, provoca que estos dispositivos sean mas cerrados al 
usuario; poco personalizables; usando tanto hardware y software de su 
propiedad.
\end{comment}
\begin{comment}
    La impresión 3D se ha consolidado como una de las tecnologías de mayor crecimiento en entornos industriales, educativos y de prototipado rápido. Lo que comenzó en laboratorios de investigación especializados se ha transformado en una herramienta omnipresente en la industria, talleres y hogares de aficionados. Esta popularidad ha generado una amplia oferta de arquitecturas comerciales que se proyectan como dispositivos comunes para el consumo masivo.

Sin embargo, esta masificación ha traído consigo un problema de diseño: la mayoría de estos equipos están destinados a un usuario no especializado que busca una experiencia de operación inmediata con mínima intervención técnica. Esto provoca que los dispositivos sean sistemas cerrados, poco personalizables y basados en hardware y software propietario.

Para una institución educativa o un laboratorio de investigación, este modelo de "caja negra" acarrea dificultades operativas graves:

Dependencia de repuestos oficiales: Cualquier avería obliga a la compra de componentes costosos y específicos del fabricante.

Dependencia administrativa: En el contexto de una universidad, depender de garantías oficiales o procesos de compra lentos a menudo condena a los equipos al desuso ante fallas menores.

Limitación del aprendizaje: Al no poder modificar el firmware o intervenir el hardware, el estudiante se limita a ser un operador de la máquina, perdiendo la oportunidad de entender la ingeniería detrás del proceso.

La transición hacia un equipo basado en tecnología de código abierto (Open Source) permite que las reparaciones y mejoras sean constantes, de bajo costo y accesibles mediante la modificación del firmware y la fabricación de piezas propias.

Por lo tanto, surge la necesidad de desarrollar una plataforma de impresión 3D abierta, modular y reparable que permita tanto el aprendizaje técnico como la reducción de costos operativos en entornos académicos.
\end{comment}